Un proton, initialement au repos, est accéléré par une différence de potentiel de \qty{\v{dV}}{kV}. Il pénètre ensuite dans une région où règne un champ magnétique uniforme de \qty{\v{B}}{mT}, perpendiculaire à sa vitesse. ($m_p = \qty{1.673e-27}{kg}$, $e = \qty{1.602e-19}{C}$)

\begin{enumerate}[label=\alph*)]
    \item Déterminez la vitesse du proton à son entrée dans le champ magnétique.
    \item Déterminez le rayon de sa trajectoire circulaire.
    \item Déterminez la période de son mouvement circulaire.
\end{enumerate}
